\src{176}
Let $\Omega_X$ (respectively $\Omega_Y$) be the sheaf of germs of regular
differential forms on $X$ (respectively $Y$). Then there are canonical
isomorphisms
\begin{equation}\tag{3.6}
  \Extsh_{\Osh_X}^{p}\bigl(\Omega_Y^{n-p},\Omega_X^{n}\bigr)=\Osh_Y,
\end{equation}
or equivalently
\begin{equation}\tag{3.6 bis}
  \Extsh_{\Osh_X}^{p}\bigl(\Osh_Y,\Omega_X^{n}\bigr)=\Omega_Y^{n-p}.
\end{equation}
Formula (3.6 bis) can serve as the definition of $\Omega_Y^{n-p}$ when $Y$ is a
singular variety. More precisely:

\result{PROPOSITION}{5} Let $Y$ be an algebraic subset of dimension $q=n-p$ in
a nonsingular algebraic variety $X$ of dimension $n$. Let $F$ be a coherent
algebraic sheaf on $X$ whose support is contained in $Y$, and let $L$ be a
locally free algebraic sheaf on $X$. Then the sheaves
$\Extsh_{\Osh_X}^{i}(F,L)$ vanish for $i<p$, while for $i=p$ there is a canonical
isomorphism
\begin{equation}\tag{3.7}
  \Extsh_{\Osh_X}^{p}(F,L)
  =\Homsh_{\Osh_X}\!\left(F,
  \Extsh_{\Osh_X}^{p}(\Osh_X/J,L)\right),
\end{equation}
where $J$ denotes an arbitrary sheaf of ideals on $X$ that annihilates $F$ and
whose zero set is $Y$. In particular, if $F$ is a coherent algebraic sheaf on
$Y$, then
\begin{equation}\tag{3.7 bis}
  \Extsh_{\Osh_X}^{p}(F,L)
  =\Homsh_{\Osh_Y}\!\left(F,\Extsh_{\Osh_X}^{p}(\Osh_Y,L)\right).
\end{equation}
Finally, still taking $F$ to be a coherent algebraic sheaf on $Y$, the sheaves
\[
  E^{i}(F)=\Extsh_{\Osh_X}^{p+i}(F,\Omega_X^{n})
\]
do not depend on the choice of embedding of the algebraic space $Y$ into a
nonsingular variety $X$.

Since the question is local, we may suppose that $X$ is affine and
$L=\Osh_X$. This reduces matters to a question in commutative algebra, indeed in
local algebra: if $A$ is a regular local ring and $M$ an $A$-module whose
support has dimension $\leq q=n-p$, one must prove that
$\operatorname{Ext}_{A}^{i}(M,A)=0$ for $i<p$, and that
\[
  \operatorname{Ext}_{A}^{p}(M,A)
  =\operatorname{Hom}_{A}\!\left(M,\operatorname{Ext}_{A}^{p}(A/J,A)\right),
\]
where $J$ is any ideal of “dimension” $\leq q$ that annihilates $M$. For the
first assertion, argue by induction on $q$: an immediate dévissage reduces to
the case $M=A/J$; replacing $J$ by a smaller ideal, using the induction
hypothesis and the exact sequence of $\operatorname{Ext}$ groups, then reduces
to the case where $J$ is generated by a “system of parameters,” as in
Proposition 4, where the result is immediate. The preceding result implies that
if $J$ is an ideal

